\documentclass[a4paper,12pt]{article}
\usepackage[utf8]{inputenc}
\usepackage[T1]{fontenc}
\usepackage[ngerman]{babel}
\usepackage{amsmath, amssymb}
\usepackage{geometry}
\usepackage{tikz}
\usetikzlibrary{shapes.geometric, arrows.meta, positioning, fit, backgrounds}
\usepackage{parskip}
\usepackage{titlesec}
\usepackage{booktabs}
\usepackage{enumitem}

\geometry{left=3cm, right=3cm, top=2.5cm, bottom=2.5cm}

\titleformat{\section}{\large\bfseries}{}{0em}{}[\titlerule]
\titleformat{\subsection}{\normalsize\bfseries}{}{0em}{}

\tikzset{
  box/.style={rectangle, draw, rounded corners=3pt, text width=4.2cm,
              align=center, minimum height=0.9cm, font=\small, inner sep=4pt},
  boxw/.style={box, text width=5.2cm},
  boxn/.style={box, text width=3.8cm},
  arrow/.style={-{Latex[length=2mm]}, thick},
}

\begin{document}

\begin{center}
  {\LARGE\bfseries Studierhinweise zu Lektion 1}\\[0.5em]
  {\large Analysis (Modul 61211)}
\end{center}

\vspace{1em}

Diese Kurseinheit dient der Einführung des Raumes $\mathbb{R}^n$ und einiger seiner
Strukturen. Die ersten beiden Abschnitte bieten einen Rückblick auf den Kurs
Mathematische Grundlagen (MG) -- wer den Stoff noch gut im Gedächtnis hat,
kann sie zunächst überfliegen, sollte aber die eingefügten Ergänzungen nicht
verpassen.

Vom dritten Abschnitt an werden neue Begriffsbildungen vorgestellt, denen sich
$\mathbb{R}^n$ als wichtiges Beispiel unterordnet: reeller Vektorraum, normierter Raum,
Umgebung eines Punktes. Besonders wichtig sind der Konvergenzbegriff und der
Vollständigkeitsbegriff (Abschnitt~1.5). Abgesehen davon, dass es sich bei den
neuen Begriffen um genaue Analogien zu schon Bekanntem aus MG handelt,
besitzen die meisten eine anschaulich-geometrische Deutung, sodass man mit
ihnen bald sicher umgehen kann.

% -----------------------------------------------------------------------
\section*{Struktur der Kurseinheit 1}

\begin{center}
\begin{tikzpicture}[node distance=0.5cm and 0.6cm]

  % Column headers
  \node[font=\bfseries\small] (hA) at (0,0)    {Rückblick auf MG};
  \node[font=\bfseries\small] (hB) at (9.5,0)  {Ergänzungen und neue Inhalte};

  % ---- Left column: MG-Rückblick ----
  \node[box, below=0.8cm of hA] (k1) {1.1 Körper $(\mathbb{R},+,\cdot)$\\Rechenregeln};
  \node[box, below=0.4cm of k1] (k2) {1.1 $(\mathbb{R},\leq)$\\linear geordneter Körper};
  \node[box, below=0.4cm of k2] (k3) {1.1 Ungleichungen,\\Betrag $|\cdot|$};
  \node[box, below=0.4cm of k3] (k4) {1.1 beschränkte Mengen,\\Funktionen, Folgen};
  \node[box, below=0.4cm of k4] (k5) {1.1 Supremum, Infimum\\Intervalle};
  \node[box, below=0.8cm of k5] (k6) {1.2 Abstand, $\varepsilon$-Umgebung\\Umgebung in $\mathbb{R}$};
  \node[box, below=0.4cm of k6] (k7) {1.2 konvergente Folge,\\Grenzwert, Regeln};
  \node[box, below=0.4cm of k7] (k8) {1.2 $\varepsilon$-$n_0$-Kriterium,\\Monotoniekriterium,\\Cauchykriterium};
  \node[box, below=0.4cm of k8] (k9) {1.2 Cauchyfolge in $\mathbb{R}$,\\Satz von Bolzano-Weierstraß};

  % ---- Right column: new content ----
  \node[box, below=0.8cm of hB] (n1) {1.1 geordnete Menge,\\geordneter Körper};
  \node[box, below=0.4cm of n1] (n2) {1.1 Cauchy-Schwarzsche,\\Minkowskische Ungleichung};
  \node[box, below=0.4cm of n2] (n3) {1.1 Räume $B(M)$,\\Supremumnorm $\|\cdot\|_\infty$};
  \node[box, below=0.4cm of n3] (n4) {1.3 $\mathbb{R}^n$ als Menge\\der $n$-Tupel; Addition,\\skalare Multiplikation};
  \node[box, below=0.4cm of n4] (n5) {1.3 reeller Vektorraum,\\Beispiele};
  \node[box, below=0.4cm of n5] (n6) {1.4 Norm $\|\cdot\|$, normierter\\Raum $(X,\|\cdot\|)$,\\induzierter Abstand};
  \node[box, below=0.4cm of n6] (n7) {1.4 Skalarprodukt,\\euklidische Norm,\\Satz des Pythagoras};
  \node[box, below=0.4cm of n7] (n8) {1.5 $\varepsilon$-Umgebung, konv.\\Folge in $(X,\|\cdot\|)$};
  \node[box, below=0.4cm of n8] (n9) {1.5 Cauchyfolge,\\Banachraum, $\mathbb{R}^n$ vollst.};
  \node[box, below=0.4cm of n9] (n10){1.5 Äquivalenz der\\Normen auf $\mathbb{R}^n$};

  % Arrows left column
  \draw[arrow] (k1) -- (k2);
  \draw[arrow] (k2) -- (k3);
  \draw[arrow] (k3) -- (k4);
  \draw[arrow] (k4) -- (k5);
  \draw[arrow] (k5) -- ++(0,-0.6) -- (k6.north);
  \draw[arrow] (k6) -- (k7);
  \draw[arrow] (k7) -- (k8);
  \draw[arrow] (k8) -- (k9);

  % Arrows right column
  \draw[arrow] (n1) -- (n2);
  \draw[arrow] (n2) -- (n3);
  \draw[arrow] (n3) -- (n4);
  \draw[arrow] (n4) -- (n5);
  \draw[arrow] (n5) -- (n6);
  \draw[arrow] (n6) -- (n7);
  \draw[arrow] (n7) -- (n8);
  \draw[arrow] (n8) -- (n9);
  \draw[arrow] (n9) -- (n10);

  % Cross arrows (left -> right)
  \draw[arrow, dashed] (k5.east) -- ++(0.3,0) |- (n3.west);
  \draw[arrow, dashed] (k5.east) -- ++(0.6,0) |- (n4.west);
  \draw[arrow, dashed] (k8.east) -- ++(0.3,0) |- (n8.west);
  \draw[arrow, dashed] (k9.east) -- ++(0.3,0) |- (n9.west);

\end{tikzpicture}
\end{center}

\newpage

% -----------------------------------------------------------------------
\section*{Zielelement 1.1 -- Rückblick und Ergänzungen: Reelle Zahlen}

\subsection*{Lerninhalte}

\begin{center}
\begin{tikzpicture}[node distance=0.5cm and 1.0cm]
  \node[box] (koer) {$(\mathbb{R},+,\cdot)$ ist ein Körper\\Körpereigenschaften};
  \node[box, right=1.0cm of koer] (ord) {geordnete Menge,\\linear geordneter Körper};
  \node[box, below=0.5cm of koer] (ungl) {Ungleichungen, Betrag $|\cdot|$\\Dreiecksungleichung};
  \node[box, below=0.5cm of ungl] (csm)  {Cauchy-Schwarzsche und\\Minkowskische Ungleichung};
  \node[box, below=0.5cm of csm]  (bschr){beschränkte Mengen, Funktionen, Folgen};
  \node[box, below=0.5cm of bschr](sup)  {Supremum, Infimum\\Existenz, Eindeutigkeit};
  \node[box, below=0.5cm of sup]  (int)  {Intervalle\\offene, halboffene, abgeschlossene};
  \node[box, below=0.5cm of int]  (ext)  {erweiterte Zahlengerade $\widehat{\mathbb{R}}$\\$\sup M, \inf M \in \widehat{\mathbb{R}}$};
  \node[box, below=0.5cm of ext]  (bm)   {Räume $\mathrm{Abb}(M,\mathbb{R})$ und $B(M)$\\Supremumnorm $\|\cdot\|_\infty$};
  \draw[arrow] (koer) -- (ungl);
  \draw[arrow] (ord.south) |- (ungl.east);
  \draw[arrow] (ungl) -- (csm);
  \draw[arrow] (csm) -- (bschr);
  \draw[arrow] (bschr) -- (sup);
  \draw[arrow] (sup) -- (int);
  \draw[arrow] (int) -- (ext);
  \draw[arrow] (ext) -- (bm);
\end{tikzpicture}
\end{center}

\subsection*{Lernziele}

Nach Durcharbeiten dieses Abschnitts sollten Sie:
\begin{itemize}
  \item die Körpereigenschaften von $(\mathbb{R},+,\cdot)$ und die Ordnungseigenschaften von
        $(\mathbb{R},\leq)$ benennen und erläutern können,
  \item Ungleichungen mit dem Betrag sicher handhaben und die Cauchy-Schwarzsche
        sowie Minkowskische Ungleichung kennen,
  \item die Begriffe beschränkte Menge, Supremum und Infimum definieren und die
        Existenz des Supremums in $\mathbb{R}$ kennen,
  \item die Intervalltypen unterscheiden und die erweiterte Zahlengerade $\widehat{\mathbb{R}}$
        beschreiben können,
  \item den Raum $B(M)$ der beschränkten Funktionen und die Supremumnorm
        $\|\cdot\|_\infty$ kennen.
\end{itemize}

\subsection*{Selbstkontrollelement 1.1}

Sehen Sie, dass die Menge $\{2^n \mid n \in \mathbb{N}\}$ nicht nach oben beschränkt
ist? -- Können Sie einen Beweis angeben?

\newpage

% -----------------------------------------------------------------------
\section*{Zielelement 1.2 -- Rückblick und Ergänzungen: Konvergenz}

\subsection*{Lerninhalte}

\begin{center}
\begin{tikzpicture}[node distance=0.5cm and 1.0cm]
  \node[box] (bet)  {Betrag, Abstand $d(x,y) = |x-y|$};
  \node[box, below=0.5cm of bet]  (umg) {$\varepsilon$-Umgebung $U_\varepsilon(a)$,\\Umgebung, Hausdorffeigenschaft};
  \node[box, below=0.5cm of umg]  (flg) {konvergente Folge,\\Grenzwert $\lim_{k\to\infty} a_k$};
  \node[box, below=0.5cm of flg]  (reg) {Rechenregeln für konvergente\\Folgen, Sandwich-Theorem};
  \node[box, right=0.8cm of flg]  (kri) {Konvergenzkriterien:\\notwendig: beschränkt\\hinreichend: Monotonie\\notw.$+$hinr.: Cauchy};
  \node[box, below=0.5cm of reg]  (tf)  {Teilfolgen, Divergenzkriterium};
  \node[box, below=0.5cm of tf]   (cau) {Cauchyfolge, Cauchykriterium};
  \node[box, below=0.5cm of cau]  (bw)  {Satz von Bolzano-Weierstraß};
  \node[box, below=0.5cm of bw]   (bdi) {bestimmt divergente Folgen,\\uneigentlicher Grenzwert};
  \draw[arrow] (bet) -- (umg);
  \draw[arrow] (umg) -- (flg);
  \draw[arrow] (flg) -- (reg);
  \draw[arrow] (kri.south) |- (reg.east);
  \draw[arrow] (reg) -- (tf);
  \draw[arrow] (tf) -- (cau);
  \draw[arrow] (cau) -- (bw);
  \draw[arrow] (bw) -- (bdi);
\end{tikzpicture}
\end{center}

\subsection*{Lernziele}

Nach Durcharbeiten dieses Abschnitts sollten Sie:
\begin{itemize}
  \item den Begriff der Umgebung mit der $\varepsilon$-Umgebung in Verbindung bringen
        und die wichtigsten Eigenschaften von Umgebungen angeben können,
  \item das $\varepsilon$-$n_0$-Kriterium, das Monotonie\-kriterium und das Cauchykriterium
        für konvergente Folgen kennen und anwenden können,
  \item mit Teilfolgen und dem Divergenzkriterium arbeiten können,
  \item den Satz von Bolzano-Weierstraß formulieren und anwenden können,
  \item bestimmt divergente Folgen und uneigentliche Grenzwerte erklären können.
\end{itemize}

\subsection*{Selbstkontrollelement 1.2}

Können Sie einige Bedingungen (notwendig und/oder hinreichend) dafür nennen,
dass eine Folge keine Cauchyfolge ist?

\newpage

% -----------------------------------------------------------------------
\section*{Zielelement 1.3 -- $\mathbb{R}^n$ als reeller Vektorraum}

\subsection*{Lerninhalte}

\begin{center}
\begin{tikzpicture}[node distance=0.5cm and 1.0cm]
  \node[box] (rn)  {Die Menge $\mathbb{R}^n$:\\senkrecht angeordnete $n$-Tupel};
  \node[box, below=0.5cm of rn]  (ops) {Addition und\\skalare Multiplikation in $\mathbb{R}^n$};
  \node[box, right=0.8cm of ops] (geo) {Veranschaulichung:\\Punkte, Ortsvektoren,\\freie Vektoren};
  \node[box, below=0.5cm of ops] (vr)  {reeller Vektorraum $(X, +, \cdot)$:\\Axiome Add$_{1-4}$, Mult$_{1-3}$};
  \node[box, below=0.5cm of vr]  (bsp) {$\mathbb{R}^n$ als reeller Vektorraum;\\weitere Beispiele:\\$\mathrm{Abb}(M,\mathbb{R})$, $B(M)$, $c$, $P_n$};
  \draw[arrow] (rn) -- (ops);
  \draw[arrow] (geo.south) |- (vr.east);
  \draw[arrow] (ops) -- (vr);
  \draw[arrow] (vr) -- (bsp);
\end{tikzpicture}
\end{center}

\subsection*{Lernziele}

Nach Durcharbeiten dieses Abschnitts sollten Sie:
\begin{itemize}
  \item den Raum $\mathbb{R}^n$ (Elemente, Addition, skalare Multiplikation) beschreiben
        und für $n = 2, 3$ geometrisch veranschaulichen können,
  \item die Axiome eines reellen Vektorraums aufzählen und überprüfen können,
  \item $\mathbb{R}^n$ als Beispiel eines reellen Vektorraums einordnen und weitere
        Beispiele (Funktionenräume, Folgenräume) kennen.
\end{itemize}

\subsection*{Selbstkontrollelement 1.3}

Sei $X$ ein reeller Vektorraum und $\emptyset \neq U \subseteq X$. Ist Ihnen klar,
dass $U$ genau dann ein Unterraum von $X$ ist, wenn für alle $\alpha, \beta \in \mathbb{R}$
und für alle $x, y \in U$ der Vektor $\alpha x + \beta y$ in $U$ liegt?

\newpage

% -----------------------------------------------------------------------
\section*{Zielelement 1.4 -- $\mathbb{R}^n$ als normierter Raum}

\subsection*{Lerninhalte}

\begin{center}
\begin{tikzpicture}[node distance=0.5cm and 1.0cm]
  \node[box] (nor)  {Norm $\|\cdot\|$ auf einem\\reellen Vektorraum $X$:\\Definitheit, Homogenität,\\Dreiecksungleichung};
  \node[box, below=0.5cm of nor]  (nrm) {normierter Raum $(X, \|\cdot\|)$,\\induzierter Abstand\\$d(x,y) := \|x-y\|$};
  \node[box, below=0.5cm of nrm]  (bsp) {Beispiele: $\|\cdot\|_\infty$ auf $B(M)$,\\$\|\cdot\|_1$ und $\|\cdot\|_\infty$ auf $C([a,b])$};
  \node[box, below=0.5cm of bsp]  (rnn) {Normen auf $\mathbb{R}^n$:\\$\|\cdot\|_1$, $\|\cdot\|_2$ (euklidisch),\\$\|\cdot\|_\infty$ (Maximum)};
  \node[box, right=0.8cm of rnn]  (ska) {Skalarprodukt $x \cdot y$,\\orthogonale Vektoren};
  \node[box, below=0.5cm of rnn]  (euk) {euklidische Norm:\\$|x| = \sqrt{x \cdot x}$};
  \node[box, below=0.5cm of euk]  (pyt) {Satz des Pythagoras};
  \draw[arrow] (nor) -- (nrm);
  \draw[arrow] (nrm) -- (bsp);
  \draw[arrow] (bsp) -- (rnn);
  \draw[arrow] (ska.south) |- (euk.east);
  \draw[arrow] (rnn) -- (euk);
  \draw[arrow] (euk) -- (pyt);
\end{tikzpicture}
\end{center}

\subsection*{Lernziele}

Nach Durcharbeiten dieses Abschnitts sollten Sie:
\begin{itemize}
  \item den Begriff des normierten (reellen Vektor-)Raums definieren und an
        Beispielen, insbesondere am Beispiel $\mathbb{R}^n$, erläutern können,
  \item den von einer Norm induzierten Abstand definieren können,
  \item das Skalarprodukt auf $\mathbb{R}^n$ mit der euklidischen Norm in Zusammenhang
        bringen und den Satz des Pythagoras formulieren können.
\end{itemize}

\subsection*{Selbstkontrollelement 1.4}

Sei $(X, \|\cdot\|)$ ein normierter Raum, $d$ der von $\|\cdot\|$ induzierte Abstand.
Sehen Sie, dass $d$ translationsinvariant ist, d.\,h.\ $d(x+a, y+a) = d(x,y)$ für alle
$x, y, a \in X$ gilt?

\newpage

% -----------------------------------------------------------------------
\section*{Zielelement 1.5 -- Konvergenz in $\mathbb{R}^n$}

\subsection*{Lerninhalte}

\begin{center}
\begin{tikzpicture}[node distance=0.5cm and 1.0cm]
  \node[box] (umg)  {$\varepsilon$-Umgebung, Umgebung\\in $(X, \|\cdot\|)$,\\Hausdorffeigenschaft};
  \node[box, below=0.5cm of umg]  (kfl) {konvergente Folge $(x_k)$,\\Grenzwert in $(X,\|\cdot\|)$,\\$\varepsilon$-$n_0$-Kriterium};
  \node[box, below=0.5cm of kfl]  (krd) {Konvergenz in $(\mathbb{R}^n, \|\cdot\|_\infty)$\\$\Leftrightarrow$ Konvergenz jeder\\Koordinatenfolge};
  \node[box, right=0.8cm of krd]  (reg) {Rechenregeln für\\konvergente Folgen\\in $(X, \|\cdot\|)$};
  \node[box, below=0.5cm of krd]  (bw)  {Satz von Bolzano-Weierstraß\\in $(\mathbb{R}^n, \|\cdot\|_\infty)$};
  \node[box, below=0.5cm of bw]   (cau) {Cauchyfolge in $(X, \|\cdot\|)$,\\Eigenschaften};
  \node[box, below=0.5cm of cau]  (ban) {vollständiger normierter Raum\\(Banachraum)};
  \node[box, below=0.5cm of ban]  (rnb) {$\mathbb{R}^n$ ist Banachraum;\\$B(M)$, $C([a,b])$ vollständig};
  \node[box, below=0.5cm of rnb]  (aeq) {Äquivalenz der Normen\\auf $\mathbb{R}^n$};
  \draw[arrow] (umg) -- (kfl);
  \draw[arrow] (kfl) -- (krd);
  \draw[arrow] (reg.south) |- (bw.east);
  \draw[arrow] (krd) -- (bw);
  \draw[arrow] (bw) -- (cau);
  \draw[arrow] (cau) -- (ban);
  \draw[arrow] (ban) -- (rnb);
  \draw[arrow] (rnb) -- (aeq);
\end{tikzpicture}
\end{center}

\subsection*{Lernziele}

Nach Durcharbeiten dieses Abschnitts sollten Sie:
\begin{itemize}
  \item die Begriffe $\varepsilon$-Umgebung und Umgebung in einem normierten Raum
        definieren und ihre wichtigsten Eigenschaften nennen können,
  \item den Begriff der konvergenten Folge und ihres Grenzwertes in einem normierten
        Raum definieren und an Beispielen erläutern können,
  \item die Konvergenz von Folgen in $(\mathbb{R}^n, \|\cdot\|)$ durch die Konvergenz der
        Koordinatenfolgen charakterisieren können,
  \item den Begriff der Cauchyfolge und den des vollständigen normierten Raumes
        (Banachraums) definieren und an Beispielen erläutern können,
  \item $\mathbb{R}^n$ als Beispiel eines Banachraums kennen und wissen, dass alle Normen
        auf $\mathbb{R}^n$ äquivalent sind.
\end{itemize}

\subsection*{Selbstkontrollelement 1.5}

Sehen Sie, dass eine Folge in $(\mathbb{R}^n, \|\cdot\|)$ mit beliebiger Norm genau
dann Cauchyfolge ist, wenn jede der $n$ Koordinatenfolgen Cauchyfolge in
$(\mathbb{R}, |\cdot|)$ ist?

\end{document}
